% GENERATED FILE. Edit canonical lessons, then run npm run generate:books.
% canonical-source-hash: fnv1a64:8044b7a6e8350580
% canonical-lessons: AR-C545-shair, AR-C545-qafiya, AR-C545-anshada, AR-C545-abdaa, AR-C545-akhraja

\chapter{Poets and directors}
\label{ch:ar-poets-and-directors}
\hlchaptermodality{545}

\begin{chapteropening}
\textbf{By the end of this chapter:} \emph{I can name a poet and a rhyme, and say that someone recited a song, created something brilliant or directed a film.}

Complete the last lesson of chapter 545: I can name a poet and a rhyme, and say that someone recited a song, created something brilliant or directed a film.
\end{chapteropening}

\section[\texorpdfstring{\ar{شاعر} (shāʿir)}{شاعر (shāʿir)}]{\ar{شاعر} (shāʿir) — a poet}

\label{lesson:AR-C545-shair}



\begin{quote}
Before the new one: say the Arabic for a reed flute, then the Arabic for a band.
\end{quote}

\subsection*{You'll want to know: \ar{شاعر}}
\textbf{\ar{شاعر}} — \emph{shāʿir} — \textquotedblleft{}a poet\textquotedblright{}.

\begin{quote}
\ar{كتب} \ar{الشاعر} \ar{قصيدة} \ar{عن} \ar{البحر}.
\end{quote}

\begin{quote}
\emph{kataba ash-shāʿir qaṣīda ʿan al-baḥr.}
\end{quote}

\begin{quote}
The poet wrote a poem about the sea.
\end{quote}

Plural: \ar{شعراء}, \emph{shuʿarāʾ}.

Shiʿr, poetry, has the same letters as shaʿr, hair.

\begin{cousinweb}
\textbf{\ar{شاعر}} (\emph{shāʿir}) is built on the native Arabic root \emph{sh-ʿ-r}, \textquotedblleft{}to know, to perceive\textquotedblright{} (Lane). Shuʿūr, a feeling, which you know, is from the same root: a poet is one who perceives.
\end{cousinweb}

\begin{grammarlens}[title={What you've built}]
One more word for describing what you lived and felt.
\end{grammarlens}

\subsection*{Your turn}
\begin{itemize}
\raggedright
  \item \emph{Say it:} \emph{shāʿir}
  \item \emph{Say it:} \emph{shāʿir}, once more
  \item \emph{Say it:} \emph{kataba ash-shāʿir qaṣīda ʿan al-baḥr.} — The poet wrote a poem about the sea.
  \item \emph{Recall:} say the Arabic for a reed flute, then the Arabic for a band, then say \emph{shāʿir} again
\end{itemize}

\begin{tcolorbox}[breakable,colback=teal!4,colframe=teal!35!black,title={Before you move on}]
What does \ar{شاعر} mean? (\textquotedblleft{}A poet\textquotedblright{}.) Say it once more.
\end{tcolorbox}
\section[\texorpdfstring{\ar{قافية} (qāfiya)}{قافية (qāfiya)}]{\ar{قافية} (qāfiya) — a rhyme}

\label{lesson:AR-C545-qafiya}



\begin{quote}
Before the new one: say the Arabic for a band, then the Arabic for a poet.
\end{quote}

\subsection*{You'll want to know: \ar{قافية}}
\textbf{\ar{قافية}} — \emph{qāfiya} — \textquotedblleft{}a rhyme\textquotedblright{}.

\begin{quote}
\ar{كل} \ar{بيت} \ar{في} \ar{القصيدة} \ar{له} \ar{نفس} \ar{القافية}.
\end{quote}

\begin{quote}
\emph{kull bayt fī al-qaṣīda lahu nafs al-qāfiya.}
\end{quote}

\begin{quote}
Every line in the poem has the same rhyme.
\end{quote}

\ar{بيت} is a line of verse as well as a house.

\begin{cousinweb}
\textbf{\ar{قافية}} (\emph{qāfiya}) is built on the native Arabic root \emph{q-f-w}, \textquotedblleft{}to follow\textquotedblright{} (Lane). The qāfiya is the end of the line that follows the same sound each time.
\end{cousinweb}

\begin{grammarlens}[title={What you've built}]
One more word for describing what you lived and felt.
\end{grammarlens}

\subsection*{Your turn}
\begin{itemize}
\raggedright
  \item \emph{Say it:} \emph{qāfiya}
  \item \emph{Say it:} \emph{qāfiya}, once more
  \item \emph{Say it:} \emph{kull bayt fī al-qaṣīda lahu nafs al-qāfiya.} — Every line in the poem has the same rhyme.
  \item \emph{Recall:} say the Arabic for a band, then the Arabic for a poet, then say \emph{qāfiya} again
\end{itemize}

\begin{tcolorbox}[breakable,colback=teal!4,colframe=teal!35!black,title={Before you move on}]
What does \ar{قافية} mean? (\textquotedblleft{}A rhyme\textquotedblright{}.) Say it once more.
\end{tcolorbox}
\section[\texorpdfstring{\ar{أنشد} (anshada)}{أنشد (anshada)}]{\ar{أنشد} (anshada) — to recite (verse); to chant}

\label{lesson:AR-C545-anshada}



\begin{quote}
Before the new one: say the Arabic for a poet, then the Arabic for a rhyme.
\end{quote}

\subsection*{You'll want to know: \ar{أنشد}}
\textbf{\ar{أنشد}} — \emph{anshada} — \textquotedblleft{}to recite (verse); to chant\textquotedblright{}.

\begin{quote}
\ar{أنشد} \ar{الطلاب} \ar{نشيد} \ar{المدرسة}.
\end{quote}

\begin{quote}
\emph{anshada aṭ-ṭullāb nashīd al-madrasa.}
\end{quote}

\begin{quote}
The students sang the school anthem.
\end{quote}

Present: \ar{ينشد}, \emph{yunshid}.

The noun \ar{نشيد} is an anthem or a chant.

\begin{cousinweb}
\textbf{\ar{أنشد}} (\emph{anshada}) is built on the native Arabic root \emph{n-sh-d}, \textquotedblleft{}to seek; to recite verse\textquotedblright{} (Lane). Form IV anshada is to recite verse aloud.
\end{cousinweb}

\begin{grammarlens}[title={What you've built}]
One more word for describing what you lived and felt.
\end{grammarlens}

\subsection*{Your turn}
\begin{itemize}
\raggedright
  \item \emph{Say it:} \emph{anshada}
  \item \emph{Say it:} \emph{anshada}, once more
  \item \emph{Say it:} \emph{anshada aṭ-ṭullāb nashīd al-madrasa.} — The students sang the school anthem.
  \item \emph{Recall:} say the Arabic for a poet, then the Arabic for a rhyme, then say \emph{anshada} again
\end{itemize}

\begin{tcolorbox}[breakable,colback=teal!4,colframe=teal!35!black,title={Before you move on}]
What does \ar{أنشد} mean? (\textquotedblleft{}To recite (verse); to chant\textquotedblright{}.) Say it once more.
\end{tcolorbox}
\section[\texorpdfstring{\ar{أبدع} (abdaʿa)}{أبدع (abdaʿa)}]{\ar{أبدع} (abdaʿa) — to create; to excel}

\label{lesson:AR-C545-abdaa}



\begin{quote}
Before the new one: say the Arabic for a rhyme, then the Arabic for to recite.
\end{quote}

\subsection*{You'll want to know: \ar{أبدع}}
\textbf{\ar{أبدع}} — \emph{abdaʿa} — \textquotedblleft{}to create; to excel\textquotedblright{}.

\begin{quote}
\ar{أبدعت} \ar{الفرقة} \ar{في} \ar{الحفلة}.
\end{quote}

\begin{quote}
\emph{abdaʿat al-firqa fī al-ḥafla.}
\end{quote}

\begin{quote}
The band was brilliant at the party.
\end{quote}

Present: \ar{يبدع}, \emph{yubdiʿ}.

\ar{أبدع} \ar{في} … is to do brilliantly at a thing.

\begin{cousinweb}
\textbf{\ar{أبدع}} (\emph{abdaʿa}) is built on the native Arabic root \emph{b-d-ʿ}, \textquotedblleft{}to originate, to make a thing new\textquotedblright{} (Lane). Mubdiʿ, creative, is its participle.
\end{cousinweb}

\begin{grammarlens}[title={What you've built}]
One more word for describing what you lived and felt.
\end{grammarlens}

\subsection*{Your turn}
\begin{itemize}
\raggedright
  \item \emph{Say it:} \emph{abdaʿa}
  \item \emph{Say it:} \emph{abdaʿa}, once more
  \item \emph{Say it:} \emph{abdaʿat al-firqa fī al-ḥafla.} — The band was brilliant at the party.
  \item \emph{Recall:} say the Arabic for a rhyme, then the Arabic for to recite, then say \emph{abdaʿa} again
\end{itemize}

\begin{tcolorbox}[breakable,colback=teal!4,colframe=teal!35!black,title={Before you move on}]
What does \ar{أبدع} mean? (\textquotedblleft{}To create; to excel\textquotedblright{}.) Say it once more.
\end{tcolorbox}
\section[\texorpdfstring{\ar{أخرج} (akhraja)}{أخرج (akhraja)}]{\ar{أخرج} (akhraja) — to direct films; to bring out}

\label{lesson:AR-C545-akhraja}



\begin{quote}
Before the new one: say the Arabic for to recite, then the Arabic for to create.
\end{quote}

\subsection*{You'll want to know: \ar{أخرج}}
\textbf{\ar{أخرج}} — \emph{akhraja} — \textquotedblleft{}to direct films; to bring out\textquotedblright{}.

\begin{quote}
\ar{أخرج} \ar{صديقي} \ar{فيلما} \ar{قصيرا}.
\end{quote}

\begin{quote}
\emph{akhraja ṣadīqī fīlman qaṣīran.}
\end{quote}

\begin{quote}
My friend directed a short film.
\end{quote}

Present: \ar{يخرج}, \emph{yukhrij}.

Its first sense is to take out; the director, \ar{مخرج}, is spelled like makhraj, an exit.

\begin{cousinweb}
\textbf{\ar{أخرج}} (\emph{akhraja}) is built on the native Arabic root \emph{kh-r-j}, \textquotedblleft{}to go out\textquotedblright{} (Lane). Kharaja, to go out, which you know, is from the same root; Form IV is to bring out.
\end{cousinweb}

\begin{grammarlens}[title={What you've built}]
One more word for describing what you lived and felt.
\end{grammarlens}

\subsection*{Your turn}
\begin{itemize}
\raggedright
  \item \emph{Say it:} \emph{akhraja}
  \item \emph{Say it:} \emph{akhraja}, once more
  \item \emph{Say it:} \emph{akhraja ṣadīqī fīlman qaṣīran.} — My friend directed a short film.
  \item \emph{Recall:} say the Arabic for to recite, then the Arabic for to create, then say \emph{akhraja} again
\end{itemize}

\begin{tcolorbox}[breakable,colback=teal!4,colframe=teal!35!black,title={Before you move on}]
What does \ar{أخرج} mean? (\textquotedblleft{}To direct films; to bring out\textquotedblright{}.) Say it once more.
\end{tcolorbox}
